% =====================================================================
%  Cronograma del semestre 2026-2 como imagen independiente.
%  Mismo diagrama que la Figura 1 de informe.tex, aislado para exportar.
%  Compilar:  pdflatex cronograma-gantt.tex
%  Exportar:  pdftoppm -png -r 300 cronograma-gantt.pdf cronograma-gantt
% =====================================================================
\documentclass[border=10pt]{standalone}

\usepackage[utf8]{inputenc}
\usepackage[T1]{fontenc}
\usepackage[spanish,es-tabla]{babel}
\usepackage{mathptmx}
\usepackage{tikz}
\usepackage{xcolor}

\definecolor{barfill}{HTML}{3B5C86}
\definecolor{hito}{HTML}{9B3A2E}
\definecolor{rejilla}{HTML}{C8CDD4}

\begin{document}
\begin{tikzpicture}[x=0.165cm, y=-0.68cm, font=\footnotesize]

% --- titulo
\node[anchor=south west, font=\normalsize\bfseries] at (-30,-1.95)
     {Cronograma de ejecución --- semestre 2026-2};

% --- rejilla semanal
\foreach \d in {0,7,14,21,28,35,42,49,56,63,70,77,84}{
  \draw[rejilla, line width=0.3pt] (\d,-0.45) -- (\d,7.35);
}
% --- limites de mes
\foreach \d in {0,30,61,88}{
  \draw[gray!70, line width=0.6pt] (\d,-0.75) -- (\d,7.35);
}
% --- etiquetas de mes
\node[anchor=south] at (15,-0.80) {\textsc{septiembre}};
\node[anchor=south] at (45.5,-0.80) {\textsc{octubre}};
\node[anchor=south] at (74.5,-0.80) {\textsc{noviembre}};

% --- eje
\draw[gray!70, line width=0.6pt] (0,-0.45) -- (88,-0.45);

% --- barras
\newcommand{\fase}[5]{%
  \node[anchor=east, font=\footnotesize] at (-2,#1) {#4};
  \fill[barfill] (#2,#1-0.19) rectangle (#3,#1+0.19);
  \node[anchor=west, gray!55!black, font=\scriptsize] at (#3+1.6,#1) {#5};
}
\fase{0.6}{0}{7}{\textbf{F0} Delimitación}{1--7 sep}
\fase{1.6}{7}{21}{\textbf{F1} Corpus}{8--21 sep}
\fase{2.6}{21}{35}{\textbf{F2} Protocolo}{22 sep--5 oct}
\fase{3.6}{35}{49}{\textbf{F3} Agente WikiLLM v1}{6--19 oct}
\fase{4.6}{49}{63}{\textbf{F4} Línea base}{20 oct--2 nov}
\fase{5.6}{63}{77}{\textbf{F5} Experimento}{3--16 nov}
\fase{6.6}{77}{88}{\textbf{F6} Cierre}{17--27 nov}

% --- hitos de control
\fill[hito] (35,2.6) circle (1.9pt);
\fill[hito] (63,4.6) circle (1.9pt);
\draw[hito, dashed, line width=0.5pt] (35,2.6) -- (35,7.15);
\draw[hito, dashed, line width=0.5pt] (63,4.6) -- (63,7.15);
\node[hito, anchor=north west, font=\scriptsize] at (35.5,7.15) {5 oct};
\node[hito, anchor=north west, font=\scriptsize] at (63.5,7.15) {2 nov};

% --- leyenda
\node[hito, anchor=north west, font=\scriptsize] at (-30,7.55)
     {\textbullet\ Puntos de corte: cierre de F2 y cierre de F4};

\end{tikzpicture}
\end{document}
